\usepackage[manualmark,headsepline,footsepline]{scrlayer-scrpage}

\usepackage[german]{babel}
\usepackage[utf8]{inputenc}
\usepackage{graphicx} %Grafik aktivieren
\usepackage{float}
\usepackage[final]{pdfpages}
\usepackage{hyperref}
\usepackage{pdfcomment}
\usepackage{svg}
\usepackage{csquotes}

\usepackage{siunitx}  
\sisetup{locale = DE,separate-uncertainty}  

\usepackage[backend=biber]{biblatex}
\addbibresource{bib/refs.bib}


%% used for pgfplots, tikzpicture
\usepackage{pgfplots}
\pgfplotsset{compat=newest}
\usetikzlibrary{plotmarks, arrows.meta, positioning, calc, decorations.pathreplacing} 

%% the following commands are needed for some matlab2tikz features
\usetikzlibrary{plotmarks}
\usetikzlibrary{arrows.meta}
\usetikzlibrary{datavisualization}
\usetikzlibrary{datavisualization.formats.functions}
\usepgfplotslibrary{patchplots}
\usepackage{array}
\usepackage{grffile} 
\usepackage{amsmath}

\usepackage{todonotes}


\usepackage{listings}
\usepackage{xcolor}

\definecolor{codebg}{HTML}{F5F5F5}
\definecolor{codeframe}{HTML}{D0D0D0}
\definecolor{codekw}{HTML}{0033B3}
\definecolor{codestr}{HTML}{067D17}
\definecolor{codecom}{HTML}{8C8C8C}

\lstdefinestyle{codeblock}{
    language=Python,
    backgroundcolor=\color{codebg},
    frame=single,
    rulecolor=\color{codeframe},
    framesep=4pt,
    framexleftmargin=4pt,
    xleftmargin=4pt,
    xrightmargin=4pt,
    basicstyle=\ttfamily\footnotesize,
    keywordstyle=\color{codekw}\bfseries,
    stringstyle=\color{codestr},
    commentstyle=\color{codecom}\itshape,
    showstringspaces=false,
    breaklines=true,
    breakatwhitespace=true,
    aboveskip=4pt,
    belowskip=4pt,
}
\lstset{style=codeblock}

\usepackage[most]{tcolorbox}


\usepackage[
	format=hang,			% Einstellung fuer Bildunterschriften
%	font={footnotesize},
	font={scriptsize},          % applies to both label and body
%	labelfont={bf,small},  % overrides for the label only ("Figure 3.4")
	textfont={scriptsize},      % overrides for the body only
	labelfont={bf},
	margin=1cm,
	aboveskip=5pt,
	position=bottom
	]{caption}


\pgfplotsset{plot coordinates/math parser=false}
\newlength\fwidth

%%%%%%%%%%%%%%%%%%%%%%%%%%%%%%%%%%%%%%%%%%%%%%%%%%%%%%%%%%%%%%%%%%%%%%%%%%%%%%%%%%%%%%%%%%%%%%%%%%%%%%%%%%%%%%%%%%%%%%%%

%Kapitel-Layout
\usepackage{lmodern} 
%\usepackage{xcolor}
\colorlet{chapter}{black!75}
\addtokomafont{chapter}{\color{chapter}}
\makeatletter% siehe De-TeX-FAQ
\renewcommand*{\chapterformat}{%
    \begingroup% damit \unitlength-Änderung lokal bleibt 
    \setlength{\unitlength}{1mm}% 
    \begin{picture}(20,20)(0,5)% 
    \setlength{\fboxsep}{0pt}% 
    \put(20,20){\line(1,0){\dimexpr 
    \textwidth-20\unitlength\relax\@gobble}}% 
    \put(0,0){\makebox(15,20)[r]{% 
    \fontsize{18\unitlength}{28\unitlength}\selectfont\thechapter 
    \kern-.04em% Ziffer in der Zeichenzelle nach rechts schieben 
    }}% 
    \end{picture} % <-- Leerzeichen ist hier beabsichtigt! 
    \newline
    \endgroup
}

\setlength\parindent{0pt}

%%%%%%%%%%%%%%%%%%%%%%%%%%%%%%%%%%%%%%%%%%%%%%%%%%%%%%%%%%%%%%%%%%%%%%%%%%%%%%%%%%%%%%%%%%%%%%%%%%%%%%%%%%%%%%%%%%%%%%%%

% Kopf- und Fußzeilen, Seitenränder etc.
% Kopf- und Fußzeilen ------------------------------------------------------
\pagestyle{scrheadings}

% Kopf- und Fußzeile auch auf Kapitelanfangsseiten -------------------------
\renewcommand*{\chapterpagestyle}{scrheadings}

% Schriftform der Kopfzeile ------------------------------------------------
\renewcommand{\headfont}{%
\normalfont
}

% Kopfzeile ----------------------------------------------------------------
\ihead{Regelungstechnisches Praktikum (RTP)}
\chead{}
\ohead{\includegraphics[width=0.3\textwidth]{figure/oth-regensburg-logo}}
%
\setlength{\headheight}{21mm} % Höhe der Kopfzeile
%\setheadsepline[text]{0.4pt} % Trennlinie unter Kopfzeile

% Fußzeile -----------------------------------------------------------------
\ifoot{Marius Off, Janis Graf}
\cfoot{\today}
\ofoot{\pagemark}


% erzeugt ein wenig mehr Platz hinter einem Punkt --------------------------
\frenchspacing 

\clubpenalty = 10000
\widowpenalty = 10000 
\displaywidowpenalty = 10000


\newcommand{\mt}[1]{\textrm{#1}}
